\subsection{Cierre del Proyecto - Evaluar, corregir y cerrar la entrega [Pendientes: 1. Confirmar fechas y responsables.]}\label{sec:spmp-10-3}
La metodología DAD contempla cierres de fase con productos y criterios de cumplimiento. En particular, la fase de Semi-transición comprende evaluación, correcciones, documentación final y entrega. Las fechas exactas para el desarrollo y cierre de estas actividades se encuentran actualmente \textbf{por confirmar} en el cronograma institucional. Las evidencias de los estándares y el catálogo de la sección \ref{sec:spmp-6-4} aportan la base documental del cierre~\cite{adaptacion2026,vfp2026,std12207,ieee730}:
\begin{enumerate}
\item \textbf{Evaluación.} Contrastar la versión final de la aplicación con las necesidades y los criterios de aceptación mediante la participación de los usuarios previstos, especialmente los bomberos (fecha: \textbf{por confirmar}).
\item \textbf{Correcciones.} Atender las observaciones y los problemas identificados durante la evaluación con los usuarios y comprobar los ajustes realizados (fecha: \textbf{por confirmar}).
\item \textbf{Documentación final.} Consolidar los manuales, la guía técnica, los registros de pruebas, las evidencias y los demás documentos correspondientes a la entrega (fecha: \textbf{por confirmar}).
\item \textbf{Entrega.} Preparar y presentar la versión final de la aplicación, la biblioteca de escenarios y la documentación correspondiente a sus destinatarios (fecha: \textbf{por confirmar}).
\end{enumerate}
La base de evaluación con usuarios se presenta en la sección \ref{sec:spmp-8-3}, y los responsables operativos del empaquetado y la presentación se coordinarán conjuntamente sin asignación de roles individuales previos.